\section{Verificación}
\label{sec:verificacion}

La verificación tiene cuatro niveles, de lo más aislado a lo más completo. Todas las
pruebas son automáticas y se ejecutan con \code{make test}, \code{make sim} y
\code{make sim-post}.

\subsection{Pruebas del assembler}

El archivo \code{asm/test\_pochoasm.py} contiene 20 pruebas, escritas con el módulo
\code{unittest} de Python (tabla~\ref{tab:pruebas-asm}). Las dos primeras filas son las
más importantes, porque comparan contra algo que no escribió el grupo.

\begin{itemize}
  \item \textbf{Referencia GNU.} El repositorio original trae tres programas de ejemplo
        (\code{blink.s}, \code{7seg.s}, \code{buttons\_leds.s}) con sus \code{.hex}
        generados por el assembler oficial de RISC-V de GNU. Se guardaron en
        \code{asm/golden/}, y el assembler del grupo produce exactamente las mismas 31
        palabras.
  \item \textbf{Ida y vuelta contra el decodificador.} Para cada formato con inmediato
        se codifica \emph{cada} valor posible y se decodifica con las fórmulas copiadas
        de \code{espino\_decoder.v} (por ejemplo,
        \code{imm\_b = \{\{19\{i[31]\}\}, i[31], i[7], i[30:25], i[11:8], 1'b0\}}). Tiene
        que volver el mismo número: 4.096 valores en I y S, 4.096 desplazamientos en B y
        1.048.576 en J. Si el assembler pusiera un bit en la posición equivocada, el
        procesador leería otro valor y esta prueba lo detectaría.
\end{itemize}

\begin{table}[H]
\centering
\caption{Pruebas del assembler (\code{make test}).}
\label{tab:pruebas-asm}
\begin{tabularx}{\linewidth}{l c X}
\toprule
Grupo & Pruebas & Qué se comprueba \\
\midrule
Referencia GNU & 1 & los 3 ejemplos del repositorio, palabra por palabra \\
Ida y vuelta & 5 & todos los inmediatos de I, S, B y J contra el decodificador; saltos hacia adelante y hacia atrás con etiquetas \\
Palabras conocidas & 9 & \code{nop}, \code{ret}, \code{ecall} y otras palabras de la especificación; todas las pseudo-instrucciones; \code{li} con 2.017 valores; \code{.equ}; \code{-D}; aviso de desplazamientos; comentarios y etiquetas \\
Validación & 4 & 21 programas erróneos, cada uno con su mensaje y su línea; reporte de todos los errores juntos \\
El juego & 1 & \code{game.s} se ensambla sin avisos y cabe en la RAM \\
\bottomrule
\end{tabularx}
\end{table}

\subsection{Prueba de instrucciones sobre el procesador}

El programa \code{sim/isa\_test.s} ejercita cada instrucción de RV32I en 22 pruebas:
operaciones de ALU con valores con y sin signo, inmediatos en los extremos de su rango,
\code{lui} y \code{auipc}, cada branch tomado y no tomado, \code{jal} y \code{jalr},
cargas y almacenamientos de palabras, medias palabras y bytes con extensión de signo,
lectura de los botones y del contador de ciclos, el registro \code{x0} y la ejecución de
un desplazamiento como suma. Cada prueba compara su resultado con el esperado y, si
falla, escribe su número en los displays.

El programa se ensambla con el assembler del grupo y se ejecuta sobre el RTL completo
del SoC en Icarus Verilog (\code{sim/isa\_tb.v}), usando los modelos de simulación de
Lattice para los bloques de RAM. Así se prueba el assembler contra el hardware real: si
una instrucción se codificara mal, el procesador ejecutaría otra cosa. Resultado:
\textbf{las 22 pruebas pasan}.

Para confirmar que la prueba detecta errores de verdad, se introdujeron dos errores a
propósito en el assembler. Al cambiar el \code{funct3} de \code{lbu} por el de
\code{lb}, falla la prueba 15. Al perder un bit del inmediato del formato S, el
programa se escribe encima de sí mismo y no termina. En ambos casos la simulación
reporta el fallo.

\subsection{Simulación del juego}

El testbench \code{sim/game\_tb.v} juega dos partidas completas contra el SoC corriendo
\code{game.s}, y comprueba cada regla del enunciado mirando solo los pines de la placa:
los cuatro LEDs y los catorce segmentos de los displays, que decodifica de vuelta a
dígitos igual que los leería una persona. Aprieta y suelta los botones con rebotes: tres
pulsos cortos antes de quedar firmes.

Simular 3 segundos reales son 75 millones de ciclos, demasiado para un simulador. Por
eso el testbench usa el mismo \code{game.s} ensamblado con constantes más cortas,
mediante la opción \code{-D} del assembler: \code{WAIT\_CYCLES}~=~20.000,
\code{TENTH\_CYCLES}~=~1.000, \code{DEBOUNCE\_CYCLES}~=~200 y
\code{BLINK\_CYCLES}~=~2.000. El programa es idéntico; solo cambian esos cuatro números.

La primera partida usa diez tiempos de reacción distintos, elegidos para ejercitar los
casos borde: un dígito, dos dígitos, y 120 décimas, que no cabe en los displays. Además,
en la ronda 3 aprieta primero un botón incorrecto, y en la ronda 7 mantiene un botón
apretado durante toda la fase de los cuatro LEDs. La tabla~\ref{tab:sim-juego} resume lo
que se comprobó.

\begin{table}[H]
\centering
\caption{Revisiones del testbench del juego (\code{make sim}). Todas pasan.}
\label{tab:sim-juego}
\begin{tabularx}{\linewidth}{X l}
\toprule
Regla & Resultado observado \\
\midrule
Al encender: displays en \code{00} y LEDs parpadeando & \code{00}, 0101 $\leftrightarrow$ 1010 \\
Los cuatro LEDs duran al menos \code{WAIT\_CYCLES} en cada ronda & sí, en las 21 rondas \\
Luego se enciende un solo LED & sí \\
Tiempo mostrado: 3, 7, 12, 5, 120, 1, 15, 4, 8, 6 décimas & \code{03 07 12 05 99 01 15 04 08 06} \\
Botón incorrecto: \code{EE} y la ronda se repite & \code{EE}, vuelven los 4 LEDs \\
El error no cuenta como ronda ni suma al promedio & la partida termina tras 10 correctas \\
Botón sostenido durante los 4 LEDs: no es error y el objetivo espera a que se suelte & sí, con 200 ciclos sin botones \\
Tras la décima: el último tiempo sigue visible \code{WAIT\_CYCLES} & \code{06} con los LEDs apagados \\
Promedio: $\lfloor 181 / 10 \rfloor = 18$ décimas & \code{18}, LEDs parpadeando \\
Una partida nueva empieza en \code{00} & sí \\
Dos partidas empezadas en instantes distintos dan secuencias distintas & \code{2 2 3 0 1 2 2 2 3 0} vs \code{0 2 1 1 2 3 0 3 3 0} \\
\bottomrule
\end{tabularx}
\end{table}

\subsection{Simulación post-síntesis}

Las simulaciones anteriores corren sobre el RTL, la descripción en Verilog. Lo que se
graba en la FPGA es otra cosa: el circuito que produce Yosys, hecho de tablas de
búsqueda, flip-flops y bloques de RAM de la iCE40, con el programa ya cargado en esos
bloques. Si la síntesis interpretara algo distinto del simulador, por ejemplo la
escritura por bytes de la RAM o su inicialización con \code{\$readmemh}, la
simulación del RTL no lo mostraría.

\code{make sim-post} sintetiza el diseño con la versión corta del juego, exporta el
circuito resultante como Verilog de celdas y le aplica el mismo testbench. El resultado
es idéntico al de la tabla~\ref{tab:sim-juego}: los mismos tiempos mostrados, el mismo
\code{EE}, el mismo promedio \code{18} y las mismas secuencias de LEDs en ambas
partidas. La simulación toma unos 7 minutos.

\subsection{Problemas encontrados}

El testbench encontró dos problemas durante el desarrollo, ambos en el propio
testbench y no en el juego. El primero: medía la duración de la fase de cuatro LEDs
desde que empezaba a mirarla, y no desde que los LEDs cambiaban, así que reportaba
fases cortas que no lo eran. El segundo: cambiaba los botones en el mismo instante del
flanco de reloj en que el SoC los lee, una carrera que hacía variar la semilla entre una
compilación y otra. Se corrigió asignando los botones con asignación no bloqueante.
